% ============================================================
%  EXAMEN T3L1 — MODELO A
%  Curso: 4to Año — Matemática
%  Duración: 60 minutos
%  Contenidos: circunferencia trigonométrica, ángulos en posición
%  estándar, signos por cuadrante, reducción al primer cuadrante,
%  funciones de ángulos opuestos, funciones recíprocas y
%  demostración de identidades trigonométricas.
%  Autor: Ing. Gabriel Astudillo
%
%  Nota: la clave de respuestas está al final, dentro de
%  \ifmwclaves ... \fi. Al compilar con compile_all.sh se genera
%  además la variante _scr (versión del estudiante, sin clave).
% ============================================================

\documentclass[12pt, letterpaper]{article}

% ── Paquetes ─────────────────────────────────────────────────
\usepackage{mathwizards-guia}
\mwguia{Examen — T3L1 (Modelo A)}{Ing. Gabriel Astudillo $\cdot$ 4to Año — Matemática}
\providecommand{\sen}{\sin}
\setlength{\emergencystretch}{2em}

% ── Circunferencia unitaria con cuadrantes ──────────────────
\newcommand{\mwcircunferenciaexamen}{%
\begin{tikzpicture}[scale=0.85, baseline=(current bounding box.center), line join=round]
  \draw[gray!55, -{Stealth}] (-1.36,0) -- (1.42,0);
  \draw[gray!55, -{Stealth}] (0,-1.36) -- (0,1.42);
  \draw[thick, mwprimario] (0,0) circle (1);
  \filldraw[mwprimario] (0,0) circle (1.2pt);
  \node[below left, font=\scriptsize] at (0,0) {$O$};
  \node[below right, font=\scriptsize, gray] at (1.34,0) {$x$};
  \node[above left, font=\scriptsize, gray] at (0,1.30) {$y$};
  \node[font=\bfseries, mwprimario] at (0.70,0.70) {I};
  \node[font=\bfseries, mwprimario] at (-0.70,0.70) {II};
  \node[font=\bfseries, mwprimario] at (-0.70,-0.70) {III};
  \node[font=\bfseries, mwprimario] at (0.70,-0.70) {IV};
  \node[font=\scriptsize, gray] at (1.08,0.15) {$1$};
  \node[font=\scriptsize, gray] at (0.15,1.06) {$1$};
\end{tikzpicture}}

\begin{document}

\mwportada{Examen T3L1 — Modelo A}{Circunferencia trigonométrica y reducción al primer cuadrante}{60 minutos $\cdot$ 100 puntos}

\vspace{4pt}
\begin{cuadroinfo}
  \small
  \textbf{Nombre:} \rule{0.55\linewidth}{0.4pt} \quad
  \textbf{Fecha:} \rule{0.15\linewidth}{0.4pt} \\[6pt]
  \textbf{Tiempo disponible:} 60 minutos \quad$\bullet$\quad
  \textbf{Puntaje total:} 100 puntos \\[4pt]
  \textbf{Instrucciones:}
  \begin{enumerate}[leftmargin=*, nosep, topsep=2pt]
    \item Resuelve en el espacio disponible y muestra tus pasos; si necesitas más, usa el reverso.
    \item En reducción, escribe el \textbf{cuadrante} y el \textbf{ángulo de referencia} antes del valor.
    \item No se permite calculadora: todos los valores pedidos son exactos.
    \item Revisa especialmente los \textbf{signos}: es el error más frecuente.
  \end{enumerate}
  {\footnotesize\textbf{Tiempo sugerido:} 8 $\cdot$ 7 $\cdot$ 15 $\cdot$ 8 $\cdot$ 10 $\cdot$ 12 min (secciones 1 a 6).\par}
\end{cuadroinfo}
\vspace{-6pt}

% ============================================================
\section{Circunferencia y posición estándar (16 puntos)}
% ============================================================

\begin{center}
\vspace{-10pt}\mwcircunferenciaexamen\vspace{-6pt}
\end{center}

\begin{enumerate}[leftmargin=*, label=\textbf{\arabic*.}, itemsep=0.4em]
  \item \textbf{(6 pts)} Responde brevemente:
        \begin{enumerate}[label=\alph*), leftmargin=1.6em, nosep]
          \item ¿Qué ecuación cumplen todos los puntos de la circunferencia unitaria?
          \item ¿Cuánto mide su radio?
          \item ¿Cuál es el signo del coseno en el IV cuadrante?
        \end{enumerate}

  \item \textbf{(4 pts)} Indica el cuadrante o el eje donde se ubica cada ángulo:
        $135^\circ$,\quad $250^\circ$,\quad $480^\circ$,\quad
        $-120^\circ$,\quad $270^\circ$,\quad $360^\circ$.

  \item \textbf{(3 pts)} Escribe \textbf{dos} ángulos coterminales de $130^\circ$:
        uno positivo y uno negativo. Luego halla el ángulo coterminal de
        $1030^\circ$ en el intervalo $[0^\circ,360^\circ)$.

  \item \textbf{(3 pts)} El punto $P\!\left(-\frac{3}{5},\frac{4}{5}\right)$
        pertenece a la circunferencia unitaria. Indica su cuadrante y calcula
        $\sen\alpha$, $\cos\alpha$ y $\tg\alpha$.
\end{enumerate}

\newpage
% ============================================================
\section{Signos de las funciones (12 puntos)}
% ============================================================

\begin{cuadroinfo}
  \small
  Determina el \textbf{signo} de cada expresión \textbf{sin calcular su valor}.
  Escribe el cuadrante (o el ángulo coterminal) que justifica tu respuesta.
  Cada ejercicio vale 3 puntos.
\end{cuadroinfo}

\begin{enumerate}[leftmargin=*, label=\textbf{\alph*)}, itemsep=1.1em]
  \item $\sen 160^\circ$
  \item $\cos 230^\circ$
  \item $\csc(-200^\circ)$
  \item $\cot 1030^\circ$
\end{enumerate}

% ============================================================
\section{Reducción al primer cuadrante (24 puntos)}
% ============================================================

\begin{cuadroinfo}
  \small
  Completa la tabla. Reduce cada ángulo al primer cuadrante y escribe el
  \textbf{valor exacto}. Cada fila vale 4 puntos: 1 por el cuadrante,
  1 por el ángulo de referencia y 2 por el valor.
\end{cuadroinfo}

\begin{center}
\renewcommand{\arraystretch}{2.5}
\setlength{\tabcolsep}{9pt}
\begin{tabular}{@{}llll@{}}
  \toprule
  \textbf{Ejercicio} & \textbf{Cuadrante} & \textbf{Ángulo de referencia} & \textbf{Valor exacto} \\
  \midrule
  a) $\sen 210^\circ$  & & & \\
  b) $\cos 135^\circ$  & & & \\
  c) $\tg 300^\circ$   & & & \\
  d) $\sen 480^\circ$  & & & \\
  e) $\cos 1290^\circ$ & & & \\
  f) $\tg 2115^\circ$  & & & \\
  \bottomrule
\end{tabular}
\end{center}

\newpage
% ============================================================
\section{Ángulos opuestos y funciones recíprocas (12 puntos)}
% ============================================================

\begin{cuadroinfo}
  \small
  Aplica la \textbf{paridad} de las funciones y las \textbf{razones recíprocas}.
  Cada ejercicio vale 3 puntos.
\end{cuadroinfo}

\begin{enumerate}[leftmargin=*, label=\textbf{\alph*)}, itemsep=1.2em]
  \item $\sen(-45^\circ)$
  \item $\cos(-120^\circ)$
  \item $\tg(-60^\circ)$
  \item $\sec 225^\circ$
\end{enumerate}

% ============================================================
\section{Expresiones combinadas (16 puntos)}
% ============================================================

\begin{cuadroinfo}
  \small
  Calcula el valor de cada expresión. Muestra la reducción de cada función
  antes de sustituir. Cada expresión vale 8 puntos.
\end{cuadroinfo}

\begin{enumerate}[leftmargin=*, label=\textbf{\arabic*.}, itemsep=1.5em]
  \item $E=\cos 120^\circ+\sen 210^\circ-\tg 315^\circ$
        \vspace{4em}

  \item $E=\dfrac{\sen(-60^\circ)\cdot\cot 225^\circ}{\cos 240^\circ}$
        \vspace{4em}
\end{enumerate}

\newpage
% ============================================================
\section{Demostración de identidades (20 puntos)}
% ============================================================

\begin{cuadroteoria}
  \small
  \textbf{Identidades trigonométricas fundamentales.} Una identidad es una
  igualdad que se cumple para \emph{todo} ángulo en el que las funciones estén
  definidas.
  \begin{center}
  \renewcommand{\arraystretch}{1.25}
  \begin{tabular}{@{}p{0.27\linewidth}p{0.65\linewidth}@{}}
    \toprule
    \textbf{Grupo} & \textbf{Identidades} \\
    \midrule
    Pitagórica & $\sen^{2}\alpha+\cos^{2}\alpha=1$ \\[2pt]
    Cociente & $\tg\alpha=\frac{\sen\alpha}{\cos\alpha}$, \quad
               $\cot\alpha=\frac{\cos\alpha}{\sen\alpha}$ \\[2pt]
    Recíprocas & $\csc\alpha=\frac{1}{\sen\alpha}$, \quad
                 $\sec\alpha=\frac{1}{\cos\alpha}$, \quad
                 $\cot\alpha=\frac{1}{\tg\alpha}$ \\[2pt]
    Derivadas & $1+\tg^{2}\alpha=\sec^{2}\alpha$, \quad
                $1+\cot^{2}\alpha=\csc^{2}\alpha$ \\[2pt]
    Paridad & $\sen(-\alpha)=-\sen\alpha$, \quad
              $\cos(-\alpha)=\cos\alpha$, \quad
              $\tg(-\alpha)=-\tg\alpha$ \\
    \bottomrule
  \end{tabular}
  \end{center}
\end{cuadroteoria}

\begin{cuadropasos}
  \small
  \textbf{Cómo demostrar una identidad:}
  \begin{enumerate}[leftmargin=1.4em, nosep, topsep=2pt]
    \item elige el lado \textbf{más complicado} y trabaja \emph{solo} con él;
    \item lleva todo a \textbf{seno y coseno} con las identidades de cociente y recíprocas;
    \item usa la \textbf{pitagórica}: cambia $\sen^{2}\alpha+\cos^{2}\alpha$ por $1$
          y $1-\sen^{2}\alpha$ por $\cos^{2}\alpha$;
    \item simplifica con álgebra (desarrolla, factoriza o cancela) hasta obtener
          exactamente el otro lado.
  \end{enumerate}
  No pases términos de un lado a otro: eso resuelve una ecuación, no demuestra
  una identidad.
\end{cuadropasos}

\textbf{Demuestra cada identidad.} Cada una vale \textbf{5 puntos}: 1 por
elegir el lado conveniente, 2 por usar bien las identidades y 2 por llegar a
la igualdad.

\begin{enumerate}[leftmargin=*, label=\textbf{\alph*)}, itemsep=3.2em]
  \item $\tg\alpha\cdot\cos\alpha=\sen\alpha$

  \item $\dfrac{1-\sen^{2}\alpha}{\cos\alpha}=\cos\alpha$

  \item $\sec^{2}\alpha-\tg^{2}\alpha=1$

  \item $\left(\sen\alpha+\cos\alpha\right)^{2}=1+2\,\sen\alpha\cos\alpha$
\end{enumerate}

% ============================================================
\ifmwclaves
\clearpage
\section*{Clave de respuestas comentada}
% ============================================================

\begin{cuadroinfo}
  \small
  \textbf{Cómo leer la clave.} Cada respuesta indica el camino, no solo el
  número. Si tu resultado difiere, revisa primero el \textbf{cuadrante} y el
  \textbf{signo}: son las dos fuentes de error más frecuentes.
\end{cuadroinfo}

\subsection*{Sección 1 — Circunferencia y posición estándar (16 pts)}
\begin{enumerate}[leftmargin=*, label=\textbf{\arabic*.}]
  \item a) $x^2+y^2=1$. \quad b) $r=1$. \quad c) Positivo (en el IV cuadrante $x>0$).
  \item $135^\circ$: II; \ $250^\circ$: III; \ $480^\circ\equiv120^\circ$: II; \
        $-120^\circ\equiv240^\circ$: III; \ $270^\circ$: eje $y$ negativo; \
        $360^\circ$: eje $x$ positivo.
  \item $130^\circ+360^\circ=490^\circ$ y $130^\circ-360^\circ=-230^\circ$.
        Además $1030^\circ=2\cdot360^\circ+310^\circ$, luego el coterminal es
        $310^\circ$ (IV cuadrante).
  \item II cuadrante (la $x$ es negativa y la $y$ positiva).
        \[
          \sen\alpha=\frac45,\qquad
          \cos\alpha=-\frac35,\qquad
          \tg\alpha=\frac{4/5}{-3/5}=-\frac43.
        \]
\end{enumerate}

\subsection*{Sección 2 — Signos (12 pts)}
\begin{center}
\renewcommand{\arraystretch}{1.6}
\setlength{\tabcolsep}{9pt}
\begin{tabular}{@{}llll@{}}
  \toprule
  \textbf{Expresión} & \textbf{Cuadrante} & \textbf{Signo} & \textbf{Justificación} \\
  \midrule
  $\sen 160^\circ$  & II  & $+$ & en II, $y>0$ \\
  $\cos 230^\circ$  & III & $-$ & en III, $x<0$ \\
  $\csc(-200^\circ)$ & II & $+$ & $-200^\circ\equiv160^\circ$; $\sen>0$ \\
  $\cot 1030^\circ$ & IV  & $-$ & $1030^\circ\equiv310^\circ$; $\tg<0$ \\
  \bottomrule
\end{tabular}
\end{center}

\subsection*{Sección 3 — Reducción (24 pts)}
\begin{center}
\renewcommand{\arraystretch}{1.7}
\setlength{\tabcolsep}{9pt}
\begin{tabular}{@{}llll@{}}
  \toprule
  \textbf{Ejercicio} & \textbf{Cuadrante} & \textbf{Ángulo de referencia} & \textbf{Valor exacto} \\
  \midrule
  a) $\sen 210^\circ$  & III & $30^\circ$ & $-\dfrac12$ \\
  b) $\cos 135^\circ$  & II  & $45^\circ$ & $-\dfrac{\sqrt{2}}{2}$ \\
  c) $\tg 300^\circ$   & IV  & $60^\circ$ & $-\sqrt{3}$ \\
  d) $\sen 480^\circ$  & II  & $60^\circ$ & $+\dfrac{\sqrt{3}}{2}$ \\
  e) $\cos 1290^\circ$ & III & $30^\circ$ & $-\dfrac{\sqrt{3}}{2}$ \\
  f) $\tg 2115^\circ$  & IV  & $45^\circ$ & $-1$ \\
  \bottomrule
\end{tabular}
\end{center}

\noindent Detalle de las reducciones con ángulos mayores que $360^\circ$:
\[
  480^\circ=360^\circ+120^\circ,\qquad
  1290^\circ=3\cdot360^\circ+210^\circ,\qquad
  2115^\circ=5\cdot360^\circ+315^\circ.
\]

\subsection*{Sección 4 — Opuestos y recíprocas (12 pts)}
\begin{enumerate}[leftmargin=*, label=\textbf{\alph*)}]
  \item $\sen(-45^\circ)=-\sen45^\circ=-\dfrac{\sqrt{2}}{2}$ \ (seno impar).
  \item $\cos(-120^\circ)=\cos120^\circ=-\dfrac12$ \ (coseno par).
  \item $\tg(-60^\circ)=-\tg60^\circ=-\sqrt{3}$ \ (tangente impar).
  \item $\cos225^\circ=-\dfrac{\sqrt{2}}{2}$, luego
        $\sec225^\circ=\dfrac{1}{-\sqrt{2}/2}=-\sqrt{2}$.
\end{enumerate}

\subsection*{Sección 5 — Expresiones (16 pts)}
\begin{enumerate}[leftmargin=*, label=\textbf{\arabic*.}]
  \item $\cos120^\circ=-\dfrac12$ (II), $\sen210^\circ=-\dfrac12$ (III) y
        $\tg315^\circ=-1$ (IV). Entonces
        \[
          E=-\frac12-\frac12-(-1)=0.
        \]
  \item $\sen(-60^\circ)=-\dfrac{\sqrt{3}}{2}$, $\cot225^\circ=1$ (III) y
        $\cos240^\circ=-\dfrac12$ (III). Entonces
        \[
          E=\frac{\left(-\frac{\sqrt{3}}{2}\right)(1)}{-\frac12}=\sqrt{3}.
        \]
\end{enumerate}

\clearpage
\subsection*{Sección 6 — Identidades (20 pts)}
\begin{enumerate}[leftmargin=*, label=\textbf{\alph*)}]
  \item Partimos del lado izquierdo y usamos la identidad de cociente:
        \[
          \tg\alpha\cdot\cos\alpha
          =\frac{\sen\alpha}{\cos\alpha}\cdot\cos\alpha
          =\sen\alpha .
        \]
        La igualdad vale para $\cos\alpha\neq0$, que es la condición para que
        exista la tangente.

  \item Del lado izquierdo, la pitagórica da $1-\sen^{2}\alpha=\cos^{2}\alpha$:
        \[
          \frac{1-\sen^{2}\alpha}{\cos\alpha}
          =\frac{\cos^{2}\alpha}{\cos\alpha}=\cos\alpha ,
        \]
        con $\cos\alpha\neq0$.

  \item Escribimos la secante y la tangente con seno y coseno:
        \[
          \sec^{2}\alpha-\tg^{2}\alpha
          =\frac{1}{\cos^{2}\alpha}-\frac{\sen^{2}\alpha}{\cos^{2}\alpha}
          =\frac{1-\sen^{2}\alpha}{\cos^{2}\alpha}
          =\frac{\cos^{2}\alpha}{\cos^{2}\alpha}=1 .
        \]

  \item Desarrollamos el cuadrado del binomio y agrupamos:
        \begin{align*}
          \left(\sen\alpha+\cos\alpha\right)^{2}
          &=\sen^{2}\alpha+2\,\sen\alpha\cos\alpha+\cos^{2}\alpha\\
          &=\underbrace{\left(\sen^{2}\alpha+\cos^{2}\alpha\right)}_{1}
            +2\,\sen\alpha\cos\alpha
           =1+2\,\sen\alpha\cos\alpha .
        \end{align*}
\end{enumerate}

\begin{cuadroadvertencia}
  \small
  \textbf{Error clásico en las demostraciones.} No se vale multiplicar en
  cruzado ni «pasar» términos al otro lado: eso transforma la identidad en una
  ecuación y hace circular la demostración. Se trabaja \textbf{un solo lado}
  hasta que coincida con el otro.
\end{cuadroadvertencia}

\begin{cuadroinfo}
  \small
  \textbf{Criterios de corrección sugeridos:}
  \begin{itemize}[leftmargin=1.2em, nosep, topsep=2pt]
    \item en las tablas, otorga el puntaje por paso (cuadrante, referencia,
          valor) y no solo por el resultado;
    \item un error de \textbf{signo} descuenta el paso del signo, no todo el
          ejercicio, salvo que el valor final sea lo único pedido;
    \item en una identidad, si el alumno opera en \textbf{ambos lados} a la
          vez, asigna como máximo 3 de los 5 puntos;
    \item premia el uso explícito de la identidad citada (pitagórica,
          cociente, recíproca o paridad): es lo que se está evaluando.
  \end{itemize}
\end{cuadroinfo}

\fi
% ============================================================
\end{document}
